\Needspace{22\baselineskip}
\subsection{Canonical roots of concrete states}
\label{sec:concrete-roots}
This milestone connects the existing search semantics to Appendix~C.
We construct the paper's rooted predicate, prove that search preserves it,
and establish concrete versions of Lemmas~23 and~26(1). These results use
our erasure-based search semantics on recursive state classes; identifying that
semantics with the paper's contextual local rewrites remains an adequacy
obligation. The proofs below do not assume the abstract specimen locality laws.

\begin{definition}[Canonical initial attachment]
\termdef{concrete-root}
For a concrete state $s$, erase its focus, attach a searching focus at the
sole input of the resulting graph, and call the resulting state $r(s)$.
Focus freedom and hole freedom of erasure supply the certificates required
for realization. Its focus is initial by the initial-realization theorem.
\end{definition}
\lean{Graph/StateRoots}{12}{14}{Linear.State.root}
\lean{Graph/StateRoots}{16}{17}{Linear.State.root_initial}

\begin{lemma}[Roots respect renaming and initiality]
Recursively isomorphic underlying graphs give recursively focus-preservingly
isomorphic roots. In particular, a focussed recursive isomorphism between
states gives one between their roots. If $s$ is already initial, $r(s)$ is
recursively focus-preservingly isomorphic to $s$.
\end{lemma}
\begin{proof}
The coordinate ``search at input zero'' is fixed by every boundary-preserving
port map. Apply recursive naturality of realization to the given underlying
isomorphism. For focussed state isomorphisms, first use erasure naturality.

For an initial state, the extracted mode is search and its extracted vertex
is an input. The state interface has exactly one input, so that vertex is
input zero. Thus the extracted typed coordinate equals the initial coordinate,
by component equality and subtype extensionality. Rewrite the general recovery
isomorphism using this coordinate equality, and embed the literal-label
isomorphism into a recursive one.
\end{proof}
\lean{Graph/StateRoots}{20}{25}{Linear.State.root_of_underlying}
\lean{Graph/StateRoots}{27}{28}{Linear.State.root_iso}
\lean{Graph/StateRoots}{31}{40}{Linear.State.coordinate_initial}
\lean{Graph/StateRoots}{43}{47}{Linear.State.root_initial_iso}

\begin{lemma}[Search preserves the underlying graph]
If $s\longrightarrow_d t$, their erased underlying graphs are recursively
isomorphic.
\end{lemma}
\begin{proof}
Let $q$ be the witnessing successor coordinate on the erased graph of $s$.
Insertion--erasure cancellation identifies the underlying graph of
$\operatorname{realize}(q)$ with the underlying graph of $s$. Invert that
isomorphism and compose it with erasure of the supplied recursive endpoint
isomorphism from the realization to $t$.
\end{proof}
\lean{Graph/StateRoots}{51}{56}{Linear.State.deepSearch_underlying}

Descend $r$ to recursive state classes using root isomorphism invariance.
Its value on a representative is its representative root class, by definition.
Every root class is initial. An initial class is fixed by $r$, using quotient
soundness and the initial-state recovery lemma. In particular $r(r(S))=r(S)$.
Search preserves $r(S)$ because it preserves the underlying graph; induction
on a search execution then proves preservation of the root along the whole run.
\lean{Graph/StateRoots}{63}{65}{Linear.DeepStateClass.root}
\lean{Graph/StateRoots}{67}{67}{Linear.DeepStateClass.root_ofState}
\lean{Graph/StateRoots}{69}{71}{Linear.DeepStateClass.root_initial}
\lean{Graph/StateRoots}{73}{75}{Linear.DeepStateClass.root_eq_of_initial}
\lean{Graph/StateRoots}{77}{78}{Linear.DeepStateClass.root_idempotent}
\lean{Graph/StateRoots}{81}{86}{Linear.DeepStateClass.search_root}
\lean{Graph/StateRoots}{88}{92}{Linear.DeepStateClass.steps_root}

\Needspace{20\baselineskip}
\subsection{Concrete rootedness and its search invariant}
\begin{definition}[Concrete rooted state]
\termdef{concrete-rootedness}
A recursive state class $S$ is rooted when there is a finite search-only
execution from $r(S)$ to $S$:
\[
 \mathsf{Rooted}(S)\quad\Longleftrightarrow\quad
 \exists n\in\mathbb N,\ r(S)\longrightarrow_d^n S.
\]
A named state is rooted when its recursive isomorphism class is rooted.
This definition depends on the operation classification used by search.
\end{definition}
\lean{Graph/StateRoots}{95}{96}{Linear.DeepStateClass.Rooted}
\lean{Graph/RootedStatePaths}{96}{97}{Linear.State.Rooted}

\begin{theorem}[Initial states are rooted; search preserves rootedness]
Every initial class is rooted, every canonical root is rooted, and search
steps and finite search executions preserve rootedness.
\end{theorem}
\begin{proof}
An initial class equals its root, so the zero-step run witnesses rootedness.
Every root is initial, giving the second claim. For a step $S\longrightarrow_d T$,
append that step to the witnessing run from $r(S)$ to $S$. Search preservation
of roots identifies $r(S)$ with $r(T)$, so this is the required run to $T$.
Induction on a finite execution gives the final claim.
\end{proof}
\lean{Graph/StateRoots}{98}{102}{Linear.DeepStateClass.initial_rooted}
\lean{Graph/StateRoots}{104}{105}{Linear.DeepStateClass.root_rooted}
\lean{Graph/StateRoots}{108}{113}{Linear.DeepStateClass.rooted_search}
\lean{Graph/StateRoots}{115}{119}{Linear.DeepStateClass.rooted_steps}

\begin{theorem}[Equivalent initial-run characterization]
A class $S$ is rooted exactly when some initial class has a finite search run
to $S$. The initial class in any such run must equal $r(S)$.
\end{theorem}
\begin{proof}
The forward direction chooses $r(S)$, which is initial. Conversely, suppose
$I$ is initial and reaches $S$. Initiality gives $I=r(I)$, and preservation
of roots along the run gives $r(I)=r(S)$. Rewrite the run's source by these
equalities to obtain the rootedness witness.
\end{proof}
\lean{Graph/ConcreteRefocusing}{10}{17}{Linear.DeepStateClass.rooted_iff_initial_run}

\Needspace{20\baselineskip}
\subsection{Operation reachability under coordinate search}
\begin{definition}[Operation-reachable vertex]
\termdef{operation-reachable}
A vertex $v$ is operation-reachable if there are an input index $i$ and an
edge list $E$ with an operation trace from $v$ backwards to input $i$.
Equivalently, reversing $E$ gives an operation path from that input to $v$.
An input is reachable using the empty trace. Each certified operation view
has an actual operation label, by its label certificate and the definition
of primitive operation labels at each depth.
\end{definition}
\lean{Graph/SearchReachability}{10}{11}{Linear.OperationReachable}
\lean{Graph/SearchReachability}{13}{14}{Linear.OperationReachable.input}
\lean{Graph/SearchReachability}{16}{18}{Linear.SearchOperation.isOperation}

\begin{lemma}[Argument and source reachability]
If the source of an operation is reachable, every term argument is reachable.
Conversely, if one of its arguments is reachable, its source is reachable.
\end{lemma}
\begin{proof}
Forwards, prepend the operation edge to the reverse trace from its source;
its certified argument membership supplies the target incidence. Backwards,
inspect a witnessing trace from the argument to an input. It cannot be empty,
since an edge cannot target an input. Its first edge equals the given
operation edge by uniqueness of incoming incidence. The operation's singleton
source then identifies the next vertex in the trace with its source. The
remaining trace proves source reachability.
\end{proof}
\lean{Graph/SearchReachability}{21}{25}{Linear.OperationReachable.argument}
\lean{Graph/SearchReachability}{28}{39}{Linear.OperationReachable.source}

\begin{theorem}[Search preserves operation reachability]
Every coordinate search rule preserves operation reachability of its vertex.
Recursive renaming preserves and reflects this property.
\end{theorem}
\begin{proof}
A variable trigger keeps the vertex. Operation entry either keeps its source,
for no term arguments, or moves to the first argument, using forward argument
reachability. Return from an argument first gives reachability of the operation
source by the converse lemma. It then either stays there on completion or
moves to the next argument by the forward lemma.

For renaming, transport the reverse operation trace and its edge list through
the recursive graph isomorphism. The boundary input is fixed. Apply the same
construction to the inverse isomorphism and cancel the port maps for reflection.
\end{proof}
\lean{Graph/SearchReachability}{42}{56}{Linear.SearchStep.operationReachable}
\lean{Graph/SearchReachability}{59}{62}{Linear.OperationReachable.rename}
\lean{Graph/SearchReachability}{64}{71}{Linear.OperationReachable.rename_iff}

\Needspace{20\baselineskip}
\subsection{Incoming paths and insertion of the focus}
\begin{definition}[Incoming incidence trace]
\termdef{incoming-trace}
An incoming trace is a reverse path with explicit vertex incidences but no
label restriction. It is either empty at a vertex or consists of an edge
whose target contains the current vertex, a chosen source vertex of that edge,
and a remaining incoming trace from that source.
\end{definition}
\lean{Graph/RootedPaths}{11}{14}{Linear.IncomingTrace}

\begin{lemma}[Incoming paths to reachable vertices]
An edge targeting an operation-reachable vertex is an operation, and each of
its sources is operation-reachable. Consequently, every incoming trace ending
at a reachable vertex is an operation trace with the same vertices and edge list.
\end{lemma}
\begin{proof}
Compare the incoming edge with the first edge of a witnessing reverse
operation trace to an input. The trace cannot be empty because inputs have no
incoming edge. Target-incidence uniqueness identifies the two edges, proving
its operation label. An operation has exactly one source, so every chosen
source is the next vertex in the witnessing trace and is reachable by its tail.
Induct on the arbitrary incoming trace and apply this argument at each step.
\end{proof}
\lean{Graph/RootedPaths}{18}{31}{Linear.OperationReachable.predecessor}
\lean{Graph/RootedPaths}{34}{41}{Linear.IncomingTrace.operations}
\lean{Graph/RootedPathObstruction}{11}{19}{Linear.OperationReachable.incoming_operation}

\begin{lemma}[A path never departs from its final vertex earlier]
Suppose an operation trace from $v$ to an input has length $k$. In any operation
trace to an input with length at most $k$, no edge can have $v$ as a source.
In particular the original trace never has $v$ as a source of one of its edges.
\end{lemma}
\begin{proof}
Induct on the shorter trace's edge list. If its first edge has $v$ as source,
its singleton-source property identifies the next vertex with $v$. The tail
and the original length-$k$ trace therefore both run from $v$ to inputs.
Uniqueness of operation traces to inputs equates their edge lists and hence
their lengths. But the tail is shorter than the trace of length at most $k$,
a contradiction. For a later edge, apply the induction hypothesis to the tail,
whose length is still at most $k$.
\end{proof}
\lean{Graph/RootedPaths}{45}{68}{Linear.operationTrace_sources_avoid}

\begin{theorem}[Lifting the operation path through focus insertion]
An operation path from an input to $v$ survives insertion of a focus at $v$:
it becomes a path from the same input to the hosted copy of $v$, using the
retained copies of its original edges.
\end{theorem}
\begin{proof}
The previous lemma says no source occurrence on this path equals $v$.
Insertion therefore maps all its source occurrences by hosting rather than
to the fresh focus target. Targets are always hosted. Induct on the reverse
trace, retaining each operation label and mapping each source and target
membership into its incidence list. The empty trace and the boundary input
are fixed. This proves a path in the actual inserted graph, not just in its
focus erasure.
\end{proof}
\lean{Graph/RootedPaths}{72}{95}{Linear.operationTrace_insert}

\Needspace{22\baselineskip}
\subsection{Rooted states have actual operation paths: Lemma 23}
\label{sec:rooted-operation-path}
Define the intermediate invariant $U(s)$ to mean that the extracted coordinate
of $s$ is operation-reachable in its erased graph. It is invariant under
focussed recursive isomorphisms: erasure naturality transports the graph,
coordinate naturality identifies the vertex, and reachability is invariant
under this transport.
\lean{Graph/RootedStatePaths}{12}{13}{Linear.State.UnderlyingReachable}
\lean{Graph/RootedStatePaths}{15}{21}{Linear.State.underlyingReachable_iso}

\begin{lemma}[Realization and search preserve the intermediate invariant]
Realizing a reachable typed coordinate satisfies $U$. Every canonical root
satisfies $U$, and recursive state search preserves $U$.
\end{lemma}
\begin{proof}
Coordinate recovery provides one isomorphism that restores both the original
graph and the chosen coordinate after realizing and erasing. Reflect operation
reachability along this isomorphism to prove the realization assertion.
The initial coordinate is an input, hence reachable, which proves the root
assertion. For a search step, coordinate search preserves reachability; apply
the realization assertion to its successor coordinate and transport the
result through its supplied recursive endpoint isomorphism.
\end{proof}
\lean{Graph/RootedStatePaths}{24}{33}{Linear.State.underlyingReachable_realize}
\lean{Graph/RootedStatePaths}{35}{36}{Linear.State.root_underlyingReachable}
\lean{Graph/RootedStatePaths}{39}{42}{Linear.State.deepSearch_underlyingReachable}

The invariant descends to recursive classes using its isomorphism invariance.
Quotient elimination proves the root and one-step preservation laws there.
Induction on the witnessing root run proves that every rooted class satisfies
$U$. This is the search-sequence induction used for the path argument in
Appendix~C, expressed through an invariant of erased coordinates.
\lean{Graph/RootedStatePaths}{66}{68}{Linear.DeepStateClass.UnderlyingReachable}
\lean{Graph/RootedStatePaths}{70}{72}{Linear.DeepStateClass.root_underlyingReachable}
\lean{Graph/RootedStatePaths}{74}{77}{Linear.DeepStateClass.search_underlyingReachable}
\lean{Graph/RootedStatePaths}{80}{89}{Linear.DeepStateClass.rooted_underlyingReachable}

\begin{theorem}[Concrete rooted-focus operation path]
If a concrete state is rooted, its actual focus source is reachable from
its unique input by an operation path. If the focus source is not an input,
this path is nonempty.
\end{theorem}
\begin{proof}
First consider any state satisfying $U$. Lift its witnessing erased operation
trace through insertion at its extracted coordinate, by the insertion theorem.
The general recovery isomorphism carries this inserted graph to the actual
state and carries the inserted focus source to its actual focus source.
Transport the lifted operation trace through that isomorphism. Boundary
preservation fixes the input. This gives operation reachability in the actual
state graph.

A rooted state satisfies $U$ by the invariant proof, so the construction applies.
The interface has one input, identifying its index with zero. Reverse the trace
list to obtain the existing forward operation-path representation. If this
list were empty, the zero-step reverse trace would identify the focus source
with the input; hence the path is nonempty whenever the source is not an input.
\end{proof}
\lean{Graph/RootedStatePaths}{45}{59}{Linear.State.focus_operationReachable}
\lean{Graph/RootedStatePaths}{100}{102}{Linear.State.rooted_focus_operationReachable}
\lean{Graph/RootedStatePaths}{106}{121}{Linear.State.rooted_focus_path}
This establishes the graph-path conclusion of Lemma~23 for the present concrete
search model. Unlike the paper's presentation, the proof does not first require
a separate uniqueness theorem for whole search runs: it uses the operation-path
uniqueness already proved and a direct search invariant.

\Needspace{20\baselineskip}
\subsection{The non-operation-path obstruction: Lemma 26(1)}
\label{sec:rooted-path-obstruction}
\begin{theorem}[Incoming-path obstruction]
Every incoming incidence trace to a rooted focus source is an operation trace.
Thus an incoming trace that is not an operation trace proves that the state
is not rooted.
\end{theorem}
\begin{proof}
The preceding rooted-focus theorem gives reachability of the actual source.
Apply the incoming-trace theorem at that vertex. The obstruction is its
contrapositive.
\end{proof}
\lean{Graph/RootedStatePaths}{124}{128}{Linear.State.rooted_incoming_operations}
\lean{Graph/RootedStatePaths}{132}{135}{Linear.State.not_rooted_of_nonoperation_path}

\begin{lemma}[The existing edge-path representation]
For any nonempty edge path whose last edge targets a reachable vertex,
every edge of the path is an operation, and every source of each edge is
reachable. Repeated edges are allowed in the supplied path datatype.
\end{lemma}
\begin{proof}
Induct on the path. For a single edge, use the incoming-operation and source
reachability lemmas at its target. For a path beginning with adjacent edges
$e,f$, apply the induction hypothesis to the tail beginning at $f$. The
adjacency witness is a source of $f$, hence reachable, and a target of $e$.
The incoming-edge lemmas then prove the required claims for $e$, while the
induction hypothesis proves them for all remaining edges.
\end{proof}
\lean{Graph/RootedPathObstruction}{23}{45}{Linear.path_operations_of_reachable}

\begin{theorem}[Concrete Lemma 26(1)]
If a nonempty graph path ends at the focus source of a rooted state, all of
its edges are operations. Therefore a path ending there and containing any
non-operation edge is a certificate that the state is not rooted.
\end{theorem}
\begin{proof}
Apply the edge-path lemma to the actual focus-source reachability supplied
by rootedness. For the obstruction, select the non-operation edge from the
path and contradict the resulting operation-label assertion.
\end{proof}
\lean{Graph/RootedPathObstruction}{52}{57}{Linear.State.rooted_path_operations}
\lean{Graph/RootedPathObstruction}{61}{65}{Linear.State.not_rooted_of_path}
This is the first part of the paper's obstruction lemma in our concrete path
and search representations. The one-way replacement argument in Lemma~26(2)
is proved separately in Section~\ref{sec:one-way-obstruction}. The general
output-closed specimen argument in Lemma~26(3) is proved for shallow
specimens in Section~\ref{sec:shallow-specimens}; its deep extension remains open.

\Needspace{20\baselineskip}
\subsection{Concrete refocusing and the remaining rewrite obligation}
\begin{definition}[Concrete refocusing]
\termdef{concrete-refocusing}
A transition relation on recursive state classes is refocusing when each
transition from a concretely rooted state leads to a concretely rooted state.
The rooted predicate is the one constructed from actual search above; it is
not the abstract specimen interface's supplied rooted predicate.
\end{definition}
\lean{Graph/ConcreteRefocusing}{21}{23}{Linear.DeepStateClass.Refocusing}

\begin{theorem}[Search refocusing and reduction to rewriting]
Search is refocusing. For any supplied rewrite relation $R$, the relation
$\mathsf{Search}\cup R$ is refocusing if and only if $R$ preserves rootedness.
For any machine on these concrete classes whose step relation is refocusing,
every execution from an initial class ends in a rooted class.
\end{theorem}
\begin{proof}
Search refocusing is the proved one-step rootedness invariant. If the union
is refocusing, apply its law to transitions from its $R$ branch. Conversely,
case-split a union transition: use search preservation in the search branch
and the supplied preservation law in the rewrite branch. Finally induct on
a machine execution, starting with the theorem that initial classes are rooted
and applying its refocusing hypothesis at each step.
\end{proof}
\lean{Graph/ConcreteRefocusing}{26}{28}{Linear.DeepStateClass.search_refocusing}
\lean{Graph/ConcreteRefocusing}{32}{40}{Linear.DeepStateClass.refocusing_with_search}
\lean{Graph/ConcreteRefocusing}{44}{53}{Linear.DeepStateClass.refocusing_execution}
The last two statements are conditional on the indicated rewrite or machine
preservation hypothesis. They do not construct substitution or compute
transitions, and do not assert that arbitrary rewrite relations preserve roots.

\Needspace{18\baselineskip}
\subsection{Rooted and non-rooted concrete examples}
The initial value state is rooted for any operation classification, by initiality.
Its successful passive search execution gives a rooted final value state by
preservation along search runs.
\lean{Graph/RootedExamples}{9}{10}{Linear.Examples.initial_value_rooted}
\lean{Graph/RootedExamples}{13}{14}{Linear.Examples.final_value_rooted}

The previously constructed interior-focus state is valid and hole-free but
is not rooted, for any operation classification. Its underlying closed chain
has a variable before the focus position and a value after it. In the actual
inserted graph, the retained variable edge targets the focus source. The
single-edge path consisting of that variable therefore ends at the source,
and its label is not an operation. The concrete Lemma~26(1) obstruction rules
out rootedness. This separates graph validity from the operational invariant
and exercises the obstruction on an actual named state.
\lean{Graph/RootedExamples}{18}{26}{Linear.Examples.interior_state_not_rooted}
